\chapter*{Résumé Exécutif}
\phantomsection
\addstarredchapter{Résumé Exécutif}
\fancyhead[LE]{\thepage} 
\fancyhead[RE]{Résumé Exécutif}
\fancyhead[RO]{\thepage} 
\fancyhead[LO]{Résumé Exécutif}
\fancyfoot[C]{\thepage}

\renewcommand{\headrulewidth}{0.1pt}

\setlength{\parindent}{0cm} % no indent
\setlength{\parskip}{9pt}
\vspace{-1cm}

% -------------------------------------------------------------

\hydrot\ établit un nouveau régime épistémique pour la modélisation hydrologique, remplaçant les expériences numériques faiblement couplées par des estimations synchronisées, traçables et reproductibles fondées sur la physique. Ce Livre Blanc présente une infrastructure numérique complète qui transforme la science hydrologique d'une pratique de simulations ad hoc en un processus formellement structuré d'estimation ancrée dans un continuum ontologique espace-état explicite.

\section*{Architecture et Fondements}

Le cadre conceptuel est conçu comme un artefact numérique persistant composé de six couches fonctionnelles :

\textbf{Couche Modèle} : héberge les modèles physiques à base de processus, les enregistrements de provenance logicielle et les supports spatio-temporels (grilles, discrétisations). Tous les composants du modèle sont suivis par Git, assurant la reproductibilité déterministe des noyaux numériques et de leurs environnements d'exécution.

\textbf{Couche Données} : contient les observations, réanalyses, forçages et paramètres internes du modèle. Par conception, cette couche reste privée par défaut, les paramétrages des modèles physiques constituant le seul composant potentiellement confidentiel de l'ensemble du cadre. Cet isolement réconcilie l'ouverture scientifique avec les contraintes opérationnelles et économiques.

\textbf{Couche Estimation} : effectue les opérations de comparaison, filtrage, inférence bayésienne et flux de calibration. L'accès à cette couche nécessite une authentification par clé privée. Toutes les opérations d'estimation sont en mode ajout seul : elles génèrent de nouveaux états versionnés plutôt que de modifier les données existantes, préservant ainsi des pistes d'audit complètes.

\textbf{Couche Analyse} : implémente les transformations temporelles et spatiales, les opérateurs d'extraction et les calculs de diagnostic. Cette couche fournit des modes d'accès variables selon la granularité et la sensibilité des opérations demandées.

\textbf{Couche Cartographique} : génère des visualisations et représentations spatiales compatibles avec les plateformes SIG standards. Les utilisateurs publics n'accèdent qu'à la sous-couche de rendu, recevant des sorties agrégées et des résumés statistiques sans exposition des données brutes.

\textbf{Registre Git Synchronisé} : maintient l'identité, la provenance et le versionnement à travers un registre immuable, en mode ajout seul, chaîné par hachage. Chaque événement de validation, opération de calibration et acte computationnel est enregistré cryptographiquement, assurant un engagement structurel envers l'audit et la non-répudiation.

\section*{Paradigme à Double Accès}

\hydrot\ met en œuvre un système rigoureux de contrôle d'accès dual qui équilibre la transparence scientifique avec la sécurité opérationnelle :

\subsection*{Accès Public (sans authentification)}

Les utilisateurs publics interagissent exclusivement à travers la sous-couche de rendu de la Couche Cartographique, accédant à :
\begin{itemize}
\item Cartes spatiales agrégées au niveau du bassin
\item Séries temporelles moyennées et statistiques saisonnières
\item Résumés statistiques de base (moyenne, médiane, percentiles)
\item Inspection des métadonnées révélant les versions de modèle et sources de données (sans valeurs brutes)
\item Capacités de visualisation et d'analyse exploratoire
\end{itemize}

\subsection*{Accès Authentifié par Clé Privée}

Les détenteurs de clés privées déverrouillent les capacités analytiques et opérationnelles complètes :
\begin{itemize}
\item Extraction d'hypercubes NumPy bruts à pleine résolution spatio-temporelle
\item Accès complet et droits de modification des paramètres du modèle
\item Exécution de simulations et modélisation prospective
\item Inférence bayésienne, quantification d'incertitude et flux d'ensemble
\item Opérations de calibration via la méthode \texttt{self\_fitting()}
\item Instanciation de SubTwin permettant l'extraction spatio-temporelle imbriquée
\item Opérations analytiques avancées incluant l'analyse dans le domaine fréquentiel et le calcul IA
\end{itemize}

\section*{Provenance Git et Traçabilité Cryptographique}

\subsection*{GitSynchroRecord : Identité par Provenance}

Chaque objet \hydrot\ porte un \texttt{GitSynchroRecord} capturant la provenance complète du dépôt : hachages de commit définissant les états exacts du code et des données, identifiants de branche et de tag, statut de l'arbre de travail (propre ou modifié), et URLs du dépôt distant. Ces éléments de métadonnées se combinent pour calculer une empreinte de version déterministe via hachage SHA-256.

Cette approche élimine la gestion d'identité symbolique en faveur d'une ontologie computationnelle : les objets sont définis non par des étiquettes arbitraires mais par leur lignage complet et vérifiable au sein de dépôts suivis par Git.

\subsection*{Registre Universel et Audit en Mode Ajout Seul}

Le Registre Git Universel implémente un registre immuable chaîné par hachage, stocké au format JSONL (JSON Lines). Chaque entrée du registre contient :
\begin{itemize}
\item Horodatage UTC de l'événement de validation
\item Empreinte Git-synchro de l'instance validée
\item Métadonnées d'acteur et d'hôte identifiant les parties responsables
\item Hachage de l'entrée précédente maintenant l'intégrité de la chaîne
\item Hachage de l'entrée courante calculé à partir de tous les champs précédents
\end{itemize}

Cette structure cryptographique garantit que toute altération d'une entrée invalide l'ensemble de la chaîne, fournissant des garanties d'intégrité structurelle. Toutes les actions privées---extraction de données, modifications de paramètres, exécutions de simulations---sont liées aux identités d'utilisateurs authentifiés et enregistrées de manière irrévocable, créant des pistes d'audit complètes et non répudiables.

\subsection*{Garantie de Reproductibilité}

La combinaison de provenances Git, d'empreintes déterministes et de registres en mode ajout seul fournit des garanties de reproductibilité absolue : tout état computationnel peut être précisément reconstruit à partir de son empreinte, toute modification génère une nouvelle version traçable, et toutes les lignées de dérivation restent cryptographiquement préservées sur des périodes de temps arbitraires.

\section*{Multi-Données et Multi-Modèles par Construction}

\hydrot\ est multi-données et multi-modèles par conception architecturale fondamentale, non par ajout de couches de plugins après-coup. Plusieurs produits d'observation, forçages et bases de données coexistent au sein de la même instance de jumeau numérique, chacun enregistré via sa propre identité et provenance Git. De même, plusieurs modèles physiques peuvent être liés au même système hydrologique à travers des objets logiciels et de simulation distincts.

Cette multiplicité native permet :
\begin{itemize}
\item Modélisation d'ensemble à travers des formulations structurellement différentes
\item Intercomparaison de jeux de données et réanalyses alternatives
\item Stratégies hybrides combinant des représentations physiques complémentaires
\item Raffinement spatial et temporel progressif sans tension architecturale
\item Moyennage bayésien de modèles directement fondé sur l'ontologie du jumeau numérique
\end{itemize}

\section*{Déploiement Opérationnel : Intégration \cwv}

\subsection*{Implémentation Proto : Série \hydrot\ Alpha}

Le déploiement opérationnel de \hydrot\ comme backend monolithique pour le plugin QGIS CaWaQS-Viz démontre la réalisation pratique de l'ontologie espace-état dans des environnements de production. La série d'implémentation alpha fournit :
\begin{itemize}
\item Interface publique simplifiée sans exigences d'authentification
\item Capacités de visualisation en lecture seule et d'analyse exploratoire
\item Traçabilité publique légère via réponses typées et persistance sérialisable
\item API d'accès aux données unifié abstrayant les complexités du format binaire
\end{itemize}

\subsection*{Stratégie d'Intégration Architecturale}

Le backend implémente quatre modèles de conception critiques :

\textbf{Modèle Singleton} : Une seule instance \texttt{HydroTwin} par session QGIS élimine les problèmes de cohérence d'état et fournit un accès cohérent aux données à travers plusieurs dialogues du plugin.

\textbf{API canonique en sept méthodes} : l'interface publique est désormais stabilisée autour de \texttt{configure}, \texttt{load}, \texttt{describe}, \texttt{extract}, \texttt{transform}, \texttt{render} et \texttt{export}, ce qui évite les aides transitoires et clarifie le contrat backend.

\textbf{Séparation des Préoccupations} : Des couches distinctes gèrent l'I/O binaire (Accès aux Données), les opérations NumPy pures (Couche Computationnelle), les statistiques et le bilan de masse (Couche Analyse), et le rendu de visualisation (Couche de Rendu).

\textbf{Indépendance QGIS} : L'abstraction \texttt{GeoDataProvider} découple HydroTwin des structures de données spécifiques à QGIS, permettant un déploiement autonome et l'intégration de frontends alternatifs.

\subsection*{Optimisation des Performances}

Les opérations d'I/O basées sur NumPy vectorisé réduisent les temps de lecture de données de plusieurs ordres de grandeur par rapport aux modèles d'accès séquentiels. La conversion différée vers les structures GeoDataFrame minimise l'empreinte mémoire. L'optimisation du format binaire avec découpage intelligent supporte des jeux de données dépassant la RAM disponible via des modèles d'accès mappés en mémoire.

\section*{Innovations Techniques}

\subsection*{Formalisme Espace-État Ontologique}

\hydrot\ ancre explicitement les simulations hydrologiques dans un formalisme espace-état ontologique où chaque instance représente un point ou une trajectoire bien défini dans un espace structuré de haute dimension intégrant variables physiques, forçages, paramètres et diagnostics. Ce formalisme permet des transformations composables, des comparaisons reproductibles et une quantification rigoureuse de l'incertitude.

\subsection*{Synchronisation Espace-État}

L'assimilation de données et l'inférence bayésienne sont réinterprétées comme des opérations de synchronisation espace-état plutôt que des corrections externes aux sorties du modèle. Ce changement conceptuel intègre les méthodes d'estimation directement dans l'ontologie du jumeau numérique, les traitant comme des opérations de première classe avec suivi complet de la provenance.

\subsection*{Morphismes Entre Sous-Espaces Ontologiques}

Les transformations d'échelle (interpolation, correction de biais, réduction d'échelle) sont formalisées comme des morphismes bien définis entre sous-espaces ontologiques. Ces morphismes préservent la structure sémantique tout en permettant une propagation rigoureuse des erreurs à travers les échelles spatiales et temporelles.

\subsection*{IA comme Objets Espace-État}

Les composants d'intelligence artificielle sont pleinement intégrés dans le cadre espace-état : jeux de données d'entraînement, architectures de réseaux, hyperparamètres et poids appris sont tous versionnés avec provenance Git. Cette intégration permet des méthodologies hybrides physique-IA avec traçabilité et reproductibilité complètes.

\section*{Flux de Travail de Calibration Multi-Étapes}

La méthode \texttt{self\_fitting()} supporte un flux de travail en 4 étapes :
\begin{enumerate}
\item Calibration dans le domaine fréquentiel pour l'identification paramétrique initiale
\item Optimisation temporelle de la recharge
\item Ajustement de référence hydrogéologique avec données piézométriques
\item Inférence bayésienne avec maillage adaptatif
\end{enumerate}

Toutes les opérations d'estimation nécessitent une authentification par clé privée et sont enregistrées de manière immuable dans le Registre Universel, garantissant des pistes d'audit complètes et une reproductibilité certifiée.

\section*{Cas d'Usage et Flux Opérationnels}

\subsection*{Intégrité Scientifique par Instanciation Authentifiée}

L'instanciation initiale nécessite une authentification par clé privée, déclenchant une validation complète : vérification de tous les hachages \texttt{GitSynchroRecord} pour le code, les données et les paramètres ; confirmation des arbres de travail propres ; calcul d'empreintes déterministes ; et certification du registre en mode ajout seul. Ce flux de travail garantit que seules les instances validées et traçables participent aux flux scientifiques.

\subsection*{Visualisation Publique et Analyse Exploratoire}

Les utilisateurs publics accèdent aux visualisations agrégées et résumés statistiques à travers la sous-couche de rendu, permettant l'exploration de données hydrologiques sans exposition des paramètres sensibles ou des données brutes. Cette capacité démocratise l'accès tout en préservant les contraintes de confidentialité opérationnelles.

\subsection*{Flux de Recherche Reproductible}

Les détenteurs de clés privées conduisent des recherches reproductibles avec traçabilité complète de la provenance : extraction de données brutes, modification de paramètres, exécution de simulations, quantification d'incertitude et opérations de calibration---tous enregistrés dans le Registre Universel avec identités d'acteurs et horodatages. Chaque acte computationnel génère un nouvel état versionné avec empreinte cryptographique unique.

\section*{Transformation Épistémique}

\hydrot\ transforme la modélisation hydrologique d'une séquence d'expériences numériques isolées en une épistémologie formellement structurée d'estimations physiquement fondées et synchronisées. Cette transformation :

\begin{itemize}
\item Élève les sorties de modèle au statut d'objets scientifiques de première classe avec identité et lignage cryptographiquement ancrés
\item Intègre nativement les méthodologies bayésiennes et IA dans l'ontologie du jumeau numérique
\item Réconcilie l'ouverture scientifique avec les contraintes opérationnelles via l'isolement des paramétrages privés
\item Fournit des garanties de reproductibilité absolue à travers la provenance Git déterministe
\item Établit un engagement structurel envers l'audit via des registres immuables chaînés par hachage
\item Permet une coexistence multi-modèles et multi-données sans tension architecturale
\end{itemize}

\section*{Fondation pour l'Hydrologie de Prochaine Génération}

Ce Livre Blanc présente des spécifications détaillées de la décomposition architecturale, des cas d'usage organisés par niveaux de contrôle d'accès, de la documentation technique des API, du statut de développement actuel incluant les optimisations de performance et les rationales de conception, et de la documentation complète de l'implémentation proto.

Ensemble, ces composants établissent \hydrot\ comme une infrastructure fondamentale pour la science hydrologique de prochaine génération : où chaque acte computationnel devient un objet reproductible, auditable et référençable au sein d'un jumeau numérique cohérent, évolutif et cryptographiquement ancré.

Cette transformation fournit une fondation à l'épreuve du futur pour la recherche scientifique, les prévisions opérationnelles, l'évaluation des impacts climatiques et la gestion intégrée des ressources en eau dans une ère exigeant transparence, reproductibilité et responsabilité computationnelle.

\newpage
